% ============================================================
%  Preamble — paketi i postavke
% ============================================================

% ---- Kodiranje i jezik -------------------------------------
% Kodiranje T1 drži dijakritike kao jedan znak. U OT1 su složeni, pa đ ispada
% iz tekstualnog sloja PDF-a i pretraživanje ne nalazi riječi koje ga imaju.
\usepackage[T1]{fontenc}          % UTF-8 je zadan od LaTeX 2018, inputenc nije potreban
\usepackage{lmodern}
% Engleski je glavni jezik dokumenta, a hrvatski uključuje
% \selectlanguage{croatian} prije sadržaja u main.tex i vrijedi do kraja rada.
\usepackage[croatian,english]{babel}
\usepackage{csquotes}             % biblatex ga traži uz babel
% Babel prevodi \appendix kao „Dodatak"; u hrvatskim radovima je uvriježeno „Prilog".
\addto\captionscroatian{\renewcommand{\appendixname}{Prilog}}

% ---- Raspored stranice -------------------------------------
\usepackage{geometry}
\geometry{
  a4paper,
  left=1.25in,
  right=1.25in,
  top=1in,
}
\usepackage[hang,flushmargin]{footmisc}
\usepackage{microtype}

% Numeriraju se dvije razine naslova (4.1), niže ne. Poglavlje je \section,
% odjeljak \subsection. Pododjeljci idu kao \subsubsection* i \paragraph,
% bez broja i bez stavke u sadržaju.
\setcounter{secnumdepth}{2}
\setcounter{tocdepth}{2}

% ---- Grafika i tablice -------------------------------------
\usepackage{graphicx}
\graphicspath{{slike/}}
\usepackage{float}
% Širine stupaca u tablicama pisane su u dijelovima \textwidth, a ne u
% centimetrima, pa prate marginu i ne prelijevaju se kad se ona promijeni.
\usepackage{booktabs}
\usepackage{multirow}
\usepackage{array}
\usepackage{longtable}
\usepackage{caption}
\usepackage{subcaption}
\usepackage{xcolor}
\usepackage{enumitem}
% Za listinge koji se ne smiju odvojiti od svojeg opisa.
\usepackage{needspace}

% ---- Dijagrami (TikZ) --------------------------------------
\usepackage{tikz}
\usetikzlibrary{shapes.geometric, arrows.meta, positioning, fit, calc}

% ---- Matematika --------------------------------------------
\usepackage{amsmath}
\usepackage{amssymb}

% ---- Listanje koda (PHP / Laravel) -------------------------
\usepackage{listings}
\definecolor{codebg}{rgb}{0.97,0.97,0.97}
\definecolor{codekw}{rgb}{0.12,0.29,0.55}
\definecolor{codecomment}{rgb}{0.40,0.55,0.40}
\definecolor{codestring}{rgb}{0.60,0.20,0.20}
\lstdefinestyle{kod}{
  backgroundcolor=\color{codebg},
  basicstyle=\ttfamily\footnotesize,
  keywordstyle=\color{codekw}\bfseries,
  commentstyle=\color{codecomment}\itshape,
  stringstyle=\color{codestring},
  numbers=left,
  numberstyle=\tiny\color{gray},
  numbersep=8pt,
  breaklines=true,
    showstringspaces=false,
  % Fiksni stupci ubacuju razmake oko znakova koje listings vidi kao zasebne
  % tokene, pa se `--without-intent-record` ispisivao kao `-- without - intent`.
  % Fleksibilni stupci uz keepspaces to uklanjaju, a monospace font i dalje
  % drži poravnanje ASCII tablica u dokaznim ispisima.
  columns=fullflexible,
  keepspaces=true,
  frame=single,
  rulecolor=\color{gray!40},
  tabsize=2,
  captionpos=b,
  extendedchars=true,
  literate=%
    {č}{{\v c}}1 {ć}{{\'c}}1 {ž}{{\v z}}1 {š}{{\v s}}1 {đ}{{\dj}}1
    {Č}{{\v C}}1 {Ć}{{\'C}}1 {Ž}{{\v Z}}1 {Š}{{\v S}}1 {Đ}{{\DJ}}1
}
\lstset{style=kod}
\lstdefinelanguage{PHP}{
  morekeywords={abstract,and,array,class,const,echo,else,elseif,extends,
    final,for,foreach,function,if,implements,interface,namespace,new,
    private,protected,public,return,static,throw,try,catch,finally,use,
    while,yield,match,fn,enum,readonly,null,true,false},
  sensitive=true,
  morecomment=[l]{//},
  morecomment=[l]{\#},
  morecomment=[s]{/*}{*/},
  morestring=[b]",
  morestring=[b]'
}

% ---- Bibliografija -----------------------------------------
\usepackage[backend=biber,style=ieee,sorting=none]{biblatex}

% Stil ieee ne poznaje tipove @legislation i @standard, pa je bez ovih dviju
% naredbi ispisivao „No driver for..." i posezao za zamjenskim oblikovanjem.
% Aliasi taj izbor čine izričitim: propisi i norme oblikuju se kao @misc, koji
% ispisuje polje howpublished (npr. „Narodne novine, br. 89/2025").
\DeclareBibliographyAlias{legislation}{misc}
\DeclareBibliographyAlias{standard}{misc}

\addbibresource{references.bib}

% Bilješka koja se lomi između stranica teško se čita, a rad ima nekoliko
% dugih doslovnih citata u bilješkama. Visoka kazna TeX-u govori da radije
% pomakne odlomak nego da bilješku prelomi.
\interfootnotelinepenalty=10000

% ---- Poveznice (učitati pri kraju) -------------------------
\usepackage[hidelinks]{hyperref}
\usepackage{cleveref}
% Poglavlje je \section, pa nazivi idu jednu razinu niže od uobičajenih.
\crefname{section}{poglavlje}{poglavlja}
\Crefname{section}{Poglavlje}{Poglavlja}
\crefname{subsection}{odjeljak}{odjeljci}
\Crefname{subsection}{Odjeljak}{Odjeljci}
\crefname{figure}{slika}{slike}
\Crefname{figure}{Slika}{Slike}
\crefname{table}{tablica}{tablice}
\Crefname{table}{Tablica}{Tablice}

% ---- Pomoćne naredbe ---------------------------------------
% Oznaka za nedovršeni dio teksta. Ispisuje se crveno i ulazi u popis
% na kraju rada (\popistodo), da se nedovršeno ne izgubi iz vida.
\newcounter{todocnt}
\newcommand{\todo}[1]{%
  \stepcounter{todocnt}%
  \textcolor{red}{\textbf{[TODO \thetodocnt: #1]}}%
  \addcontentsline{tdo}{section}{\protect\numberline{\thetodocnt}#1}%
}
\makeatletter
\newcommand{\popistodo}{%
  \clearpage
  \section*{Popis nedovršenog}%
  \markboth{Popis nedovršenog}{}%
  \@starttoc{tdo}%
}
\makeatother

% Strani termin pri prvom spominjanju: \strani{prekidač}{circuit breaker}
\newcommand{\strani}[2]{#1 (\emph{#2})}

% Kratica pri prvoj uporabi: \kratica{puni naziv}{engleski izvornik}{KRATICA}
\newcommand{\kratica}[3]{#1 (engl.\ \emph{#2}, skraćeno #3)}
% Kratica bez ustaljenog hrvatskog naziva: \kraticaen{KRATICA}{engleski izvornik}
\newcommand{\kraticaen}[2]{#1 (engl.\ \emph{#2})}
% Domaća kratica, bez engleskog izvornika: \kraticahr{puni naziv}{KRATICA}
\newcommand{\kraticahr}[2]{#1 (skraćeno #2)}
